\documentclass[10pt,a4paper]{amsart}
\usepackage[margin=19mm]{geometry}
\usepackage{amsmath,amssymb,mathtools}
\usepackage[scr=rsfs,cal=euler]{mathalfa}
\usepackage{xfrac}
\usepackage[unicode,hidelinks,bookmarks=false]{hyperref}
\newcommand{\ie}{i.e.\ }
\newcommand{\cf}{cf.\ }
\newcommand{\resp}{respectively\ }
\newcommand{\Subsection}{\subsection}
\providecommand{\nobreakdash}{\nobreak\hbox{-}}
\setlength{\emergencystretch}{2em}
\multlinegap0pt
\usepackage{tikz}
\newtheorem{theorem}{Theorem}
\newtheorem{lemma}{Lemma}
\newtheorem{proposition}{Proposition}
\usepackage{cleveref}
\crefname{theorem}{Theorem}{Theorems}
\crefname{lemma}{Lemma}{Lemmas}
\crefname{proposition}{Proposition}{Propositions}
\title{An Upper Bound on the Linear Tur\'{a}n Number of $k$-Crowns}
\author{Rajat Adak}
\date{Source: arXiv:2604.10467v1 (CC BY 4.0)}
\begin{document}
\maketitle

\noindent\emph{Source and licence.} An Upper Bound on the Linear Tur\'{a}n Number of $k$-Crowns. Version arXiv:2604.10467v1 (CC BY 4.0), distributed by arXiv under the Creative Commons Attribution 4.0 International licence, http://creativecommons.org/licenses/by/4.0/, as recorded in the arXiv licence metadata for this exact version and shown on the abstract page https://arxiv.org/abs/2604.10467v1. Full licence notice: \texttt{licenses/arxiv-2604.10467-CC-BY-4.0.txt}.\par\smallskip

\noindent\textit{Editorial scope.} Counted unit: the article's proposition weighted (source0.tex lines 220--225). The article's complete original proof of it is delivered with it (source0.tex lines 227--254, one \texttt{proof} environment of 28 lines). No mathematics is written, repaired or re-derived: the statement and the proof are the article's own lines, in the article's own notation, and no retained line is substituted or re-flowed.  Retained uncounted as the source-local context the proof needs: lines 92--96; lines 154--156.  Nothing was omitted from any retained span, so the package carries no omission record.  No retained line contains a citation, so the wrapper carries no bibliography and no bibliography entry.  No retained line contains a figure environment, an image-inclusion command or drawing code, so no image is delivered and the package declares no asset dependency.  Editorial formatting only, in addition: an AMS \texttt{amsart} preamble that restates the article's own declarations and macros used by the retained lines, namely the environments \texttt{theorem}, \texttt{lemma}, \texttt{proposition} on separate counters, together with the standard packages those lines need.  The retained environments print in this document as Theorem 1, Lemma 1, Proposition 1 in order of appearance, whereas the article prints the same items with the numbering of its own full document.  The complete native source of the article as distributed in the arXiv e-print for this exact version is bundled unchanged as \texttt{sources/198276/source0.tex}.
\begin{theorem}\label{main}
    For $3\leq k \leq r$,
    $$ex_r^{\mathrm{lin}}(n,C^r_{1,k}) \leq \frac{\big((k-1)(r-1)+1\big)(n-s)}{r},$$
    where $s$ is the number of vertices with degree at least $(k-1)(r-1) +2$. 
\end{theorem}

\begin{lemma}\label{lem2}
    Let $e = \{u_1,u_2,\dots u_r\}$ be any arbitrary edge in $H$ with $d(u_1) \geq d(u_2) \geq\dots \geq d(u_r)$. Then there exists $i \in [k]$ such that $d(u_i) \leq (k-i)(r-1) +1$.
\end{lemma}

\begin{proposition}\label{weighted}
Let $H$ be a linear $C_{1,k}^r$-free $r$-graph on $n$ vertices, and let $s$ denote the number of vertices of degree at least $(k-1)(r-1)+2$. Then
\[
\sum_{e\in E(H)} \frac{1}{k(e)(r-1)+1} \le \frac{n-s}{r}.
\]
\end{proposition}

\begin{proof}
We use the same notation and the same double-counting argument as in the proof of \Cref{main}. Let $I$ be the indicator function defined there. Then
$\sum_{v\in V(H)} I(v)=n-s$.

Fix an edge $e=\{u_1,\dots,u_r\}\in E(H)$ with $d(u_1)\ge d(u_2)\ge \cdots \ge d(u_r)$,
and set $t:=k(e)+1$. Since $e$ is not the base edge of any copy of $C_{1,t}^r$, the argument of Lemma \ref{lem2} yields an index $i\in [t]$ such that
\[
d(u_i)\le (t-i)(r-1)+1.
\]
Hence, exactly as in the proof of \Cref{main},
\[
\sum_{v\in e}\frac{I(v)}{d(v)}
\ge \frac{r-i+1}{(t-i)(r-1)+1}
\ge \frac{r}{(t-1)(r-1)+1}
= \frac{r}{k(e)(r-1)+1}.
\]
Summing over all edges and reversing the order of summation, we obtain
\[
r\sum_{e\in E(H)}\frac{1}{k(e)(r-1)+1}
\le
\sum_{e\in E(H)}\sum_{v\in e}\frac{I(v)}{d(v)}
=
\sum_{v\in V(H)} I(v)
=
n-s.
\]
This proves the proposition.
\end{proof}

\end{document}
